% discrete_diff_geometry.tex
% Section 4 -- discrete differential geometry on the bipartite
% octahedral lattice.  Discrete exterior calculus (forms, d, codifferential),
% discrete connections and curvature, integration by parts.

\label{sec:discrete_diff_geometry}

\placeholder{Section body -- write during Phase 3.  Develop the
discrete exterior calculus on the lattice: discrete forms, discrete
exterior derivative $d$, discrete codifferential, Stokes' theorem
analogue, discrete connections, and discrete curvature.}

\subsection{Discrete Forms}
\label{subsec:dec_forms}

\placeholder{Define discrete $k$-forms on the lattice as functions on
$k$-cells (vertices, edges, plaquettes, \ldots).  Address the
bipartite colour structure: how forms transform under colour-class
exchange.}

\subsection{Discrete Exterior Derivative}
\label{subsec:dec_d}

\placeholder{Define the discrete exterior derivative $d \colon
\Omega^k \to \Omega^{k+1}$ on the lattice.  Prove $d \circ d = 0$.}

\subsection{Discrete Codifferential and Inner Product}
\label{subsec:dec_codifferential}

\placeholder{Inner product on discrete forms (using the discrete
measure of Section~\ref{sec:discrete_measure_theory}).  Define the
codifferential $d^\ast$ as the formal adjoint of $d$.  State the
discrete Hodge decomposition.}

\subsection{Discrete Stokes' Theorem and Integration by Parts}
\label{subsec:dec_stokes}

\placeholder{Lattice-Stokes identity and the integration-by-parts
formula that follows.  This is the load-bearing identity for the
function-space theory in Section~\ref{sec:function_spaces}.}

\subsection{Discrete Connections and Curvature}
\label{subsec:dec_connections}

\placeholder{Discrete connections on the lattice (gauge / parallel-
transport viewpoint); discrete curvature as the failure of the
plaquette holonomy to commute.  Cross-reference Paper~II's gauge
derivation as a downstream consumer.}
